% 07 · HDBSCAN* — density at every scale
\section{HDBSCAN$^*$: A HIERARCHY OVER EVERY DENSITY}
\label{sec:hdbscan}

\deckslides{slides 22--26 --- mutual reachability, the tree, and the long-lived branches.}

If the problem with DBSCAN is one global density threshold, the obvious repair
is to use all of them at once and let the data choose. That is what
HDBSCAN$^*$ does \citep{Campello:13, McInnes:17}: it builds a single
hierarchy of density levels, scores each candidate cluster by how long it
survives as the threshold moves, and returns the candidates that persist.
It is the clusterer attached to t-SNE and UMAP in our pipeline, and --- in a
different guise --- it is inside EVoC.

\subsection{From points to a density-aware metric}

The construction starts by making the metric itself density-aware. For each
point, compute the core distance $\kappa_k(\mathbf{x})$ of \S\ref{sec:knn}:
the distance to its $k$-th nearest neighbour. Then replace the original
distance between two points with the \emph{mutual reachability distance},

\begin{equation}
d_{\rm mreach}(\mathbf{a},\mathbf{b}) = \max\left\{\kappa_k(\mathbf{a}),\
\kappa_k(\mathbf{b}),\ d(\mathbf{a},\mathbf{b})\right\}.
\label{eq:mreach}
\end{equation}

Read it as a floor: two points in a dense region are pushed apart to at least
the local density scale, so a sparse point can never be ``close'' to anything.
Under this metric, a dense region becomes a compact object even when its
members are far apart in the original geometry --- which is exactly the
property DBSCAN lacked.

\begin{marginnote}[]
\entry{Mutual reachability}{The distance used by HDBSCAN$^*$: the maximum of
the two core distances and the true distance \textbf{(Equation~\ref{eq:mreach})}.
It makes dense regions self-similar across scales.}
\entry{Condensed tree}{The hierarchy after small clusters have been
progressively collapsed into their parents as the density threshold rises. Its
branches are the candidate clusters, weighted by lifetime (persistence).}
\end{marginnote}

\begin{figure}[htbp]\centering
  \begin{tabular}{@{}c@{\hskip8pt}c@{}}
    \fig[0.40]{mutual_reachability.png} &
    \fig[0.50]{knn_core_distance.png}
  \end{tabular}
  \caption{The ingredients of \textbf{Equation~\ref{eq:mreach}}, drawn in the
  deck. Left: two points, each ringed by its own core distance, with the
  straight-line distance between them; the mutual reachability distance is the
  largest of the three quantities. Right: the core distance itself --- the
  radius that contains exactly $k$ neighbours --- which is the local density
  estimate of \S\ref{sec:knn} written as a length.}
  \label{fig:mreach}
\end{figure}

\subsection{Building the hierarchy}

With the metric fixed, the rest is single-linkage clustering at every
threshold simultaneously, computed efficiently:

\begin{enumerate}
\item Build the minimum spanning tree (MST) of the mutual-reachability graph.
  One tree contains every single-linkage merge at every density level: the
  largest edge weight on the path between two points is the density at which
  they merge.
\item Sort the tree's edges by weight and cut them in that order, producing a
  dendrogram of nested partitions --- the density hierarchy
  (\textbf{Figure~\ref{fig:hdbscantree}}).
\item Condense: a cluster that splits into pieces, one of which is smaller
  than \texttt{min\_cluster\_size} and short-lived, is treated as noise rather
  than as a new cluster. This is what makes the tree robust against
  single-linkage's notorious chaining.
\item Score every remaining branch by \emph{stability}: the sum, over its
  members, of the difference between the density at which the member leaves
  the cluster and the density at which the cluster was born
  (\textbf{Equation~\ref{eq:stability}}). Selecting the set of branches with
  the largest total stability --- choosing a parent over its children when the
  parent scores higher --- gives the final labels.
\end{enumerate}

\begin{equation}
S(C_j) = \sum_{\mathbf{x} \in C_j} \left( \lambda_{\mathbf{x}} -
\lambda_{\rm birth}(C_j) \right),
\qquad \lambda = 1/d_{\rm mreach} .
\label{eq:stability}
\end{equation}

\begin{figure}[htbp]\centering
  \fig[0.72]{hdbscan_tree.png}
  \caption{The density hierarchy as a single-linkage dendrogram with a
  horizontal density cut. Everything above the line is one cluster; everything
  below it has split. HDBSCAN$^*$ does not commit to the line --- it scores
  every branch by how far it survives above its own birth level and keeps the
  branches with the best scores.}
  \label{fig:hdbscantree}
\end{figure}

\begin{figure}[htbp]\centering
  \fig[0.94]{hdbscan_film.png}
  \caption{The same data at four density levels, from the deck's animation:
  clusters appear as the threshold drops, some grow, and some merge into each
  other. HDBSCAN$^*$ sees the whole sequence at once and keeps the objects that
  live long enough to be structure rather than coincidence.}
  \label{fig:hdbscanfim}
\end{figure}

\begin{figure}[htbp]\centering
  \fig[0.62]{hdbscan_sweep.png}
  \caption{One density level of the sweep. Density-based clustering is not a
  fixed partition of the data but a family of partitions indexed by scale; the
  algorithm's job is to pick from the family without being told the scale in
  advance.}
  \label{fig:hdbscansweep}
\end{figure}

\subsection{What the knobs mean}

HDBSCAN$^*$ replaces DBSCAN's two knobs with two different ones, and both are
easier to reason about:

\begin{itemize}
\item \texttt{min\_cluster\_size} is a persistence floor, not a density
  threshold. It says: a structure smaller than this many points is not
  evidence of a cluster, so collapse it into its parent rather than report it.
  Our default is 5 --- deliberately low, because several of our real clusters
  are small --- and \S\ref{sec:benchmark} reports the effect of raising it to
  10, which on the DR17 data was the single largest improvement in recall.
\item \texttt{min\_samples} sets the $k$ inside the core distance of
  \textbf{Equation~\ref{eq:mreach}}: how many neighbours define the local
  density scale. Larger values give a smoother, more conservative density
  estimate --- the usual bias--variance trade of \S\ref{sec:knn}.
\item The metric is a third, silent knob: the whole construction presumes that
  a distance is meaningful. Cosine and Euclidean give different hierarchies;
  that is why \S\ref{sec:data} normalises rows, and why the fairness argument
  for comparing methods rests on the metric being shared.
\end{itemize}

\subsection{Why it fits our problem, and where it still fails}

The hierarchy is a natural fit for a field of 25 clusters of very different
sizes: it can, in principle, keep a 6-member cluster at one density and a
230-member cluster at another. In practice two things still bite.

First, \texttt{min\_cluster\_size} is a \emph{single} size floor. It cannot be
right for the Pleiades with 23 candidates and for M~67 with 230 members at the
same time --- one of them will be over-merged or the other fragmented. The
next chapter removes this knob entirely.

Second, and more fundamentally, the algorithm inherits the geometry it is
given. On our real abundances the residual failure is not the hierarchy but
the space: the field is a single dense blob that contains the clusters, and no
choice of \texttt{min\_cluster\_size} separates an object from its own
background. That is what motivates the embeddings of
\S\ref{sec:tsne}--\S\ref{sec:evoc}: change the geometry, then let HDBSCAN$^*$
work as designed.

\begin{issues}[EXERCISES]
\begin{enumerate}
\item For the four points $\mathbf{a},\mathbf{b},\mathbf{c},\mathbf{d}$ of a
  square grid, compute $d_{\rm mreach}$ with $k=2$ and verify that it is a
  valid metric. Which of the metric axioms does it break if you omit the
  $\kappa$ terms and use $\max\{d(\mathbf{a},\mathbf{b}),\ \text{const}\}$?
\item Compute the stability score of \textbf{Equation~\ref{eq:stability}} for
  a two-branch tree in which one branch contains 100 points that survive two
  decades of $\lambda$ and the other contains 10 points that survive five.
  Which does the algorithm prefer, and does that match your intuition about
  which is a better cluster?
\item Run HDBSCAN$^*$ on the DR19 member-plus-field matrix with row
  normalisation on and off. Record the number of clusters and the largest
  cluster fraction in each case, and identify which setting produces the
  ``one blob'' failure of \S\ref{sec:dbscan}.
\end{enumerate}
\end{issues}
